Napriek uvedeným výhodám proces čítania DNA pomocou nanopórov prináša špecifické výzvy.
Pri sekvenovaní DNA cez nanopór sa DNA molekula prechádza cez nanopór po jednotlivých nukleotidoch,
pričom elektrický signál identifikuje bázy.
%todo: doplniť odkaz na sekciu o nanopórovom sekvenovaní
Podrobnejšie vysvetlenie tejto technológie bude uvedené v sekcii \ref{sec:nanopore_sequencing}.     
 Tento proces je však náchylný na chyby,
keďže DNA sa môže pohybovať rýchlo alebo nepravidelne, čo vedie k nesprávnemu
určeniu sekvencie nukleotidov. Tieto chyby v čítaní majú priamy vplyv na integritu
uložených dát a vyžadujú sofistikované metódy kódovania a korekcie chýb.
S pokračujúcim vývojom technológií nanopórového sekvenovania sa
presnosť postupne zvyšuje, čo robí z DNA perspektívnu alternatívu k tradičným metódam ukladania dát.
